```latex
\documentclass[11pt]{article}

\usepackage[a4paper,margin=1in]{geometry}
\usepackage{amsmath,amssymb,amsthm,mathtools}
\usepackage{bm}
\usepackage[colorlinks=true,linkcolor=blue,citecolor=blue,urlcolor=blue]{hyperref}

\numberwithin{equation}{section}

\newtheorem{theorem}{Theorem}[section]
\newtheorem{lemma}[theorem]{Lemma}
\newtheorem{proposition}[theorem]{Proposition}
\newtheorem{corollary}[theorem]{Corollary}

\theoremstyle{remark}
\newtheorem{remark}[theorem]{Remark}

\DeclareMathOperator{\divg}{div}
\DeclareMathOperator{\supp}{supp}
\DeclareMathOperator{\diam}{diam}
\DeclareMathOperator{\KL}{KL}

\newcommand{\R}{\mathbb R}
\newcommand{\E}{\mathbb E}
\newcommand{\Prob}{\mathbb P}
\newcommand{\bone}{\mathbf 1}
\newcommand{\cF}{\mathcal F}
\newcommand{\cG}{\mathcal G}
\newcommand{\cV}{\mathcal V}
\newcommand{\cH}{\mathcal H}
\newcommand{\eps}{\varepsilon}
\newcommand{\lesssim}{\lesssim}
\newcommand{\gtrsim}{\gtrsim}
\newcommand{\asymp}{\asymp}
\newcommand{\norm}[1]{\left\|#1\right\|}
\newcommand{\inner}[2]{\left\langle #1,#2\right\rangle}

\begin{document}

\title{Local minimax estimation of the conductivity in Darcy's problem near a constant}
\author{}
\date{}
\maketitle

\begin{abstract}
We determine, up to logarithmic factors, the local minimax rate for estimating
the conductivity in the divergence-form Darcy equation on a Euclidean ball from
interior Gaussian white noise observations of the solution.  For conductivity
perturbations of fixed small H\"older--Zygmund size about the constant function
$\bone$, the linearized forward operator is microlocally smoothing of order one
rather than two.  The resulting mean squared error rate is
$n^{-2\alpha/(2\alpha+d+2)}$, modulo polylogarithmic factors.  The smaller value
of the perturbation radius is used to make the quadratic part of the nonlinear
forward map negligible at the minimax resolution.  The proof combines elliptic
Schauder estimates, a microlocal Weyl law for the linearized Darcy operator,
Gaussian sequence-space arguments, and a finite-dimensional sieve estimator.
\end{abstract}

\section{Result}

Let $D=\{x\in\R^d:|x|<1\}$, $d\ge 2$, and write $C^\alpha(D)$ for the
H\"older--Zygmund space $B^\alpha_{\infty,\infty}(D)$.  For $r\in(0,1)$, define
\[
\cF(\alpha,r)
=
\left\{f=1+h:\ h\in C^\alpha(D),\quad
       \|h\|_{C^\alpha(D)}\le r\right\}.
\]
Every $f\in\cF(\alpha,r)$ satisfies
\[
1-r\le f\le 1+r,
\]
and hence is uniformly elliptic.  For such $f$, let $u_f\in H^1_0(D)$ be the
unique weak solution of
\begin{equation}\label{eq:Darcy}
\divg(f\nabla u_f)=1\quad\text{in }D,
\qquad
u_f=0\quad\text{on }\partial D.
\end{equation}

We use the standard Hilbert-space formulation of Gaussian white noise.  Let
$\mathbb W$ be the canonical isonormal Gaussian process on $L^2(D)$:
\[
\E[\mathbb W(\varphi)]=0,\qquad
\E[\mathbb W(\varphi)\mathbb W(\psi)]
=\inner{\varphi}{\psi}_{L^2(D)}.
\]
The observation is
\begin{equation}\label{eq:white-noise}
Y_n(\varphi)
=
\inner{u_f}{\varphi}_{L^2(D)}
+
\frac1{\sqrt n}\mathbb W(\varphi),
\qquad \varphi\in L^2(D).
\]
Equivalently, in every orthonormal basis $(e_j)_{j\ge1}$ of $L^2(D)$,
\[
Y_{n,j}=\inner{u_f}{e_j}_{L^2(D)}
+\frac1{\sqrt n}\xi_j,
\qquad
\xi_j\stackrel{\mathrm{i.i.d.}}{\sim}N(0,1).
\]

The main result is the following.

\begin{theorem}[Local minimax rate]\label{thm:main}
Let $d\ge2$ and $\alpha>1$.  There exist $r_0=r_0(\alpha,d)\in(0,1)$,
$C,c>0$, and $q\ge0$, depending only on $\alpha$ and $d$, such that for every
$0<r\le r_0$ and all sufficiently large $n$,
\begin{equation}\label{eq:main-rate}
c\,n^{-\frac{2\alpha}{2\alpha+d+2}}(\log n)^{-q}
\le
\inf_{\widehat f_n}
\sup_{f_0\in\cF(\alpha,r)}
\E_{f_0}\!\left[
\|\widehat f_n-f_0\|_{L^2(D)}^2
\right]
\le
C\,n^{-\frac{2\alpha}{2\alpha+d+2}}(\log n)^q .
\end{equation}
Thus, up to logarithmic factors,
\[
\boxed{\displaystyle
\inf_{\widehat f_n}
\sup_{f_0\in\cF(\alpha,r)}
\E_{f_0}\!\left[
\|\widehat f_n-f_0\|_{L^2(D)}^2
\right]
\asymp
n^{-\frac{2\alpha}{2\alpha+d+2}}.}
\]
\end{theorem}

\begin{remark}[Microlocal order of the local inverse problem]
The forward solution at $f=\bone$ is
\[
u_\bone(x)=\frac{|x|^2-1}{2d},\qquad
\nabla u_\bone(x)=\frac{x}{d}.
\]
The derivative at $\bone$ is
\[
A h=-\Delta_D^{-1}\Bigl(\frac1d\divg(xh)\Bigr),
\]
where $-\Delta_D^{-1}$ is the inverse Dirichlet Laplacian.  Formally,
$\divg(x\cdot)$ is first order and $-\Delta_D^{-1}$ is smoothing of order two,
so $A$ is generically smoothing of order one.  This explains the additional
$2$ in the denominator $2\alpha+d+2$, relative to direct nonparametric
regression.  The radial symbol $x\cdot\xi/|\xi|^2$ degenerates at $x=0$ and
in directions tangential to $x$, but the corresponding phase-space sets have
too small a volume to change the minimax exponent.

On substantially larger classes of coefficients, estimators based only on a
global $2$-ill-posed stability estimate can have the slower rate
$n^{-2(\alpha-1)/(2\alpha+d+2)}$ in squared loss; see, for example,
\cite{Nickl2020,GlobalPoly2026,PDERegressionMAP}.  Theorem~\ref{thm:main}
concerns the smaller local class centered at the explicitly computable constant
conductivity, where the order-one structure is sharp.
\end{remark}

\section{Elliptic regularity and the local expansion}

Let
\begin{equation}\label{eq:u1}
u_\bone(x)=\frac{|x|^2-1}{2d},\qquad a(x)=\nabla u_\bone(x)=\frac{x}{d}.
\end{equation}
For $h\in C^\alpha(D)$ with $\|h\|_{C^\alpha}\le r<1$, write
\begin{equation}\label{eq:Fdef}
F(h)=u_{1+h}-u_\bone.
\end{equation}
Subtracting \eqref{eq:Darcy} with $f=1+h$ from the equation for $u_\bone$
gives
\begin{equation}\label{eq:Feq}
-\divg\bigl((1+h)\nabla F(h)\bigr)=\divg(h a),
\qquad
F(h)|_{\partial D}=0.
\end{equation}

For $g\in C^\alpha(D)$ uniformly close to $1$, let
\[
\mathcal L_g w=-\divg(g\nabla w)
\]
with homogeneous Dirichlet boundary conditions, and let
\[
B_g=\mathcal L_g^{-1}.
\]
The Fr\'echet derivative of the forward map at $1+k$ is
\begin{equation}\label{eq:A-f-def}
A_{1+k}\delta
=
B_{1+k}\bigl(\divg(\delta\nabla u_{1+k})\bigr).
\end{equation}
At $k=0$ this becomes
\begin{equation}\label{eq:A-def}
A\delta
=
A_\bone\delta
=
(-\Delta_D)^{-1}\Bigl(\frac1d\divg(x\delta)\Bigr),
\end{equation}
where $(-\Delta_D)^{-1}$ solves $-\Delta w=g$ with $w|_{\partial D}=0$.

\begin{lemma}[Uniform elliptic estimates]\label{lem:elliptic}
There is $r_0\in(0,1/2)$ such that for all $r\le r_0$ and all
$k\in C^\alpha(D)$ with $\|k\|_{C^\alpha}\le r$:
\begin{enumerate}
\item The operator $B_{1+k}$ is bounded from $H^{-1}(D)$ to $H^1_0(D)$ with
operator norm bounded uniformly in $k$.
\item For every $\delta\in L^2(D)$,
\begin{equation}\label{eq:energy-A}
\|\nabla A_{1+k}\delta\|_{L^2(D)}
\le
C\|\delta\|_{L^2(D)}.
\end{equation}
\item For every $0<\beta<\alpha$,
\begin{equation}\label{eq:u-reg}
\|u_{1+k}-u_\bone\|_{C^{1+\beta}(D)}
\le
C_\beta\|k\|_{C^\beta(D)}.
\end{equation}
In particular, $\nabla u_{1+k}$ is bounded uniformly and, for fixed
$\beta\in(0,\alpha)$,
\[
\|\nabla u_{1+k}-a\|_{C^\beta(D)}\le C_\beta r.
\]
\end{enumerate}
All constants depend only on $\alpha,\beta,d$ after $r_0$ is fixed.
\end{lemma}

\begin{proof}
Uniform ellipticity follows from $1-r_0\le1+k\le1+r_0$.  The Lax--Milgram
lemma gives the first assertion.

For \eqref{eq:energy-A}, put $v=A_{1+k}\delta$ and
$q_k=\nabla u_{1+k}$.  Then
\[
-\divg((1+k)\nabla v)=\divg(\delta q_k).
\]
Testing by $v$ and using the divergence theorem yields
\[
\int_D(1+k)|\nabla v|^2\,dx
=
-\int_D\delta q_k\cdot\nabla v\,dx.
\]
Consequently
\[
(1-r)\|\nabla v\|_2^2
\le
\|\delta\|_2\|q_k\|_\infty\|\nabla v\|_2,
\]
which gives \eqref{eq:energy-A} once \eqref{eq:u-reg} supplies a uniform bound
for $q_k$.

For \eqref{eq:u-reg}, \eqref{eq:Feq} has divergence-form right-hand side
$\divg(k a)$ with $k\in C^\beta$, and uniformly elliptic $C^\beta$
coefficients.  Schauder estimates for divergence-form elliptic equations give
\[
\|F(k)\|_{C^{1+\beta}}
\le
C_\beta\|k a\|_{C^\beta}
\le
C_\beta\|k\|_{C^\beta}.
\]
The boundary of the unit ball is smooth, so the global Schauder estimate
applies.
\end{proof}

\begin{lemma}[Exact quadratic identity]\label{lem:quadratic}
Let $h,k\in C^\alpha(D)$ satisfy $\|h\|_{C^\alpha},\|k\|_{C^\alpha}\le r\le
r_0$, and set $\delta=h-k$.  Then
\begin{equation}\label{eq:quadratic-H1}
\left\|
F(h)-F(k)-A_{1+k}\delta
\right\|_{H^1_0(D)}
\le
C\|\delta\|_{L^\infty(D)}
   \|\delta\|_{L^2(D)}.
\end{equation}
\end{lemma}

\begin{proof}
Let
\[
z=F(h)-F(k),\qquad q_k=\nabla u_{1+k}.
\]
By \eqref{eq:Feq},
\begin{equation}\label{eq:z-eq}
-\divg((1+h)\nabla z)=\divg(\delta q_k).
\end{equation}
Let $z_0=A_{1+k}\delta$, so that
\begin{equation}\label{eq:z0-eq}
-\divg((1+k)\nabla z_0)=\divg(\delta q_k).
\end{equation}
Writing $z=z_0+w$ and subtracting \eqref{eq:z0-eq} from \eqref{eq:z-eq},
we obtain
\[
-\divg((1+h)\nabla w)=\divg(\delta\nabla z_0).
\]
Testing by $w$ gives
\[
(1-r)\|\nabla w\|_2^2
\le
\|\delta\|_\infty\|\nabla z_0\|_2\|\nabla w\|_2.
\]
Using \eqref{eq:energy-A} for $z_0$,
\[
\|\nabla z_0\|_2\le C\|\delta\|_2.
\]
Poincar\'e's inequality now gives \eqref{eq:quadratic-H1}.
\end{proof}

The estimate \eqref{eq:quadratic-H1} is especially useful at high frequency.
At frequency scale $m$, an $L^2$-normalized perturbation in a
$C^\alpha$-bounded family has supremum at most of order $m^{-\alpha}$.
Therefore the quadratic term is smaller than the order-one linear response
$m^{-1}$ precisely when $\alpha>1$.

\section{The microlocal structure of the derivative}

The symbol of $A=A_\bone$ is, modulo a lower-order term,
\begin{equation}\label{eq:symbol}
p(x,\xi)
=
\frac{i\,x\cdot\xi}{d|\xi|^2}.
\end{equation}
Indeed, $\divg(x\,\cdot\,)$ has principal symbol $i x\cdot\xi$, while
$(-\Delta_D)^{-1}$ has principal symbol $|\xi|^{-2}$.  Thus $A$ is a classical
pseudodifferential operator of order $-1$ away from the degenerate set
\[
\{(x,\xi):x\cdot\xi=0\}\cup\{x=0\}.
\]
The same remains true, uniformly for $f\in\cF(\alpha,r)$ with $r$ small,
because
\[
p_f(x,\xi)
=
\frac{i\,\nabla u_f(x)\cdot\xi}{f(x)|\xi|^2}
\]
and $\nabla u_f$ is close to $x/d$ in H\"older norm by
Lemma~\ref{lem:elliptic}.

The following proposition contains the microlocal information needed below.
It is a standard consequence of the semiclassical calculus and the Weyl law
for operators with symbols of variable-sign order; see
\cite{HormanderIII,BirmanSolomyak,Zworski}.  We include a detailed sketch at
the end of this section.

\begin{proposition}[Good microlocal subspaces]\label{prop:microlocal}
There are $r_0>0$, $c_0,C_0>0$, and, for every dyadic integer $m=2^J$,
a finite-dimensional subspace $V_m\subset C^\infty(\overline D)$ with the
following properties, uniformly in
$\|k\|_{C^\alpha}\le r\le r_0$.

\begin{enumerate}
\item \textbf{Dimension.}
\begin{equation}\label{eq:dim}
c_0m^d\le \dim V_m\le C_0m^d.
\end{equation}

\item \textbf{Stable approximation.}
There is a linear operator $P_m:C^\alpha(D)\to V_m$ such that
\begin{equation}\label{eq:P-stable}
\|P_mh\|_{C^\alpha}\le C_0\|h\|_{C^\alpha},
\end{equation}
\begin{equation}\label{eq:P-L2}
\|h-P_mh\|_{L^2}
\le
C_0m^{-\alpha}\|h\|_{C^\alpha},
\end{equation}
and
\begin{equation}\label{eq:P-forward}
\left\|A_{1+k}(h-P_mh)\right\|_{L^2}
\le
C_0m^{-\alpha-1}\|h\|_{C^\alpha}.
\end{equation}

\item \textbf{Restricted order-one lower bound.}
For every $q\in V_m$,
\begin{equation}\label{eq:restricted-lower}
\|A_{1+k}q\|_{L^2}
\ge
c_0m^{-1}\|q\|_{L^2}.
\end{equation}

\item \textbf{Dyadic product scale.}
There are uniformly bounded projectors $\Pi_j$ associated with $V_m$ such that
for $q\in V_m$,
\begin{equation}\label{eq:dyadic-A}
\|A_{1+k}q\|_{L^2}^2
\asymp
\sum_{j\le J}\omega_{j,m}^2\|\Pi_jq\|_2^2,
\qquad
\omega_{j,m}\ge c_0m^{-1},
\end{equation}
and
\begin{equation}\label{eq:dyadic-sup}
\|\Pi_jq\|_{L^\infty}
\le
C_02^{-j\alpha}\|q\|_{C^\alpha}.
\end{equation}

\item \textbf{Nonlinear local stability.}
If $r\le r_0$ and $h,k\in V_m$ satisfy
$\|h\|_{C^\alpha},\|k\|_{C^\alpha}\le r$, then
\begin{equation}\label{eq:nonlinear-restricted}
\left\|
F(h)-F(k)
\right\|_{L^2}
\ge
c_0m^{-1}\|h-k\|_{L^2}.
\end{equation}
\end{enumerate}
\end{proposition}

\begin{remark}
The crucial point is \eqref{eq:restricted-lower}: on the minimax resolution
space, the derivative is only one-smoothing.  No uniform lower bound of order
$m^{-1}$ holds on the whole scale-$m$ Laplacian frequency space, because modes
whose wave vector is tangential to $\nabla u_\bone$ are smoothed by order two.
Those modes occupy an asymptotically negligible phase-space fraction and are
excluded from $V_m$ at the top scales.
\end{remark}

\begin{proof}[Proof of Proposition~\ref{prop:microlocal}]
Let $J=\log_2m$.  Choose a standard dyadic semiclassical partition of unity in
$T^*D$.  At level $j$, coherent states are indexed by cells of volume
$2^{-j}$ in $x$ and angular cells on $\mathbb S^{d-1}$ in the direction of
$\xi$; their number is $O(2^{jd})$.

For $j\le J/2$, retain all coherent states.  Their total number is
$O(2^{Jd/2})=o(m^d)$.  Since any first-order divergence composed with the
inverse Dirichlet Laplacian smooths by at most two derivatives, the worst
singular scale on these levels is at least
\[
2^{-2j}\ge 2^{-J}=m^{-1}.
\]

For $J/2<j\le J$, retain only coherent states whose center $(x,\xi)$ satisfies
\begin{equation}\label{eq:good-region}
|x|\ge 2\rho,
\qquad
|x\cdot\xi|\ge c|x||\xi|,
\end{equation}
for fixed small $\rho,c>0$.  On this region,
\[
|p(x,\xi)|\ge c'|\xi|^{-1}.
\]
The same estimate remains valid with $x/d$ replaced by $\nabla u_{1+k}$,
because Lemma~\ref{lem:elliptic} gives
\[
\nabla u_{1+k}=\frac{x}{d}+O_{C^\beta}(r).
\]
The retained phase-space volume at each level is a fixed positive proportion
of the full volume; hence the total number of retained states is comparable to
$m^d$.  This proves \eqref{eq:dim}.

Wave-packet frames associated with this partition satisfy the usual
H\"older--Zygmund coefficient characterization.  Therefore discarding the
excluded states only from levels $j>J/2$ produces an $L^2$ error bounded by
\[
\left(\sum_{j>J/2}2^{-2j\alpha}\right)^{1/2}
\le Cm^{-\alpha}.
\]
This gives \eqref{eq:P-L2}.  The retained high-frequency states are good
order-one states; the excluded states are either of frequency at least $m$ or
are more strongly smoothing.  Consequently the forward error has the extra
factor $m^{-1}$ in \eqref{eq:P-forward}.

The lower bound \eqref{eq:restricted-lower} follows from the elliptic
principal symbol \eqref{eq:symbol} on the retained high-frequency region and
from the order-two lower bound $2^{-2j}\ge m^{-1}$ on the retained low
frequencies $j\le J/2$.  Asymptotic almost-orthogonality of wave packets with
separate dyadic frequency scales gives \eqref{eq:dyadic-A}.

Finally, \eqref{eq:nonlinear-restricted} follows from
Lemma~\ref{lem:quadratic} and Bony paralinearization.  If $\delta=h-k$,
then by Lemma~\ref{lem:quadratic},
\[
F(h)-F(k)-A_{1+k}\delta
\]
is controlled by products of the form $\Pi_j\delta\,\nabla A_{1+k}\Pi_j\delta$.
Using \eqref{eq:dyadic-sup},
\[
\|\Pi_j\delta\|_\infty
\le C2^{-j\alpha}r,
\]
while
\[
\|\nabla A_{1+k}\Pi_j\delta\|_2
\le C\|\Pi_j\delta\|_2
\]
by \eqref{eq:energy-A}.  Hence the quadratic remainder is bounded by
\[
Cr\sum_{j\le J}2^{-j\alpha}\|\Pi_j\delta\|_2.
\]
By Cauchy--Schwarz and \eqref{eq:dyadic-A}, this is at most
\[
Cr\left(\sum_{j\le J}
2^{2j(1-\alpha)}\omega_{j,m}^{-2}\omega_{j,m}^2
\|\Pi_j\delta\|_2^2\right)^{1/2}
\le
Cr\|A_{1+k}\delta\|_2,
\]
because $\omega_{j,m}^{-1}2^{-j}\le C$ on the retained scales and $\alpha>1$.
Choosing $r_0$ small enough makes this remainder smaller than half the linear
term.  Together with \eqref{eq:restricted-lower}, this proves
\eqref{eq:nonlinear-restricted}.
\end{proof}

\begin{lemma}[Weyl law for the linearized operator]\label{lem:Weyl}
Let $s_1(A)\ge s_2(A)\ge\cdots$ be the singular values of $A$.  Then
\begin{equation}\label{eq:Weyl}
\#\{j:s_j(A)>t\}\asymp t^{-d},
\qquad 0<t<t_0.
\end{equation}
Equivalently, $s_j(A)\asymp j^{-1/d}$.
\end{lemma}

\begin{proof}
For a classical pseudodifferential operator with principal symbol
\eqref{eq:symbol}, the singular-value counting function is given, to leading
order, by the phase-space volume
\[
N(t)
\sim
(2\pi)^{-d}
\operatorname{Vol}\left\{
(x,\xi)\in T^*D:
\frac{|x\cdot\xi|}{d|\xi|^2}>t
\right\}.
\]
See \cite{BirmanSolomyak,HormanderIII}.  For fixed $x\ne0$, write
$\xi=\rho\theta$ with $\theta\in\mathbb S^{d-1}$.  The condition is
\[
\rho<\frac{|x||\cos\angle(x,\theta)|}{dt}.
\]
Integrating first over $\theta$ and then over $\rho$ yields
\[
N(t)
\asymp
\int_D |x|^d t^{-d}\,dx
\asymp t^{-d}.
\]
The degenerate sets $x=0$ and $x\cdot\xi=0$ are hypersurfaces in phase space
and do not alter this power law.
\end{proof}

\section{The linear Gaussian inverse problem}

Before treating the nonlinear model, consider
\begin{equation}\label{eq:linear-model}
Y=Ah+\frac1{\sqrt n}\mathbb W,
\qquad
\|h\|_{C^\alpha}\le r.
\end{equation}

\begin{theorem}[Linear minimax rate]\label{thm:linear}
For $\alpha>1$ and sufficiently small $r>0$,
\[
\inf_{\widehat h_n}\sup_{\|h_0\|_{C^\alpha}\le r}
\E\|\widehat h_n-h_0\|_2^2
\asymp
n^{-\frac{2\alpha}{2\alpha+d+2}}
\]
up to polylogarithmic factors.
\end{theorem}

\begin{proof}
Choose
\begin{equation}\label{eq:balance}
m\asymp n^{\frac1{2\alpha+d+2}}.
\end{equation}

\emph{Upper bound.}
Let $\cH_m=V_m\cap\{\|h\|_{C^\alpha}\le r\}$, and define
\[
\widehat h_n
=
\arg\min_{h\in\cH_m}
\left\{-2Y(Ah)+\|Ah\|_2^2\right\}.
\]
Write $h_m=P_mh_0$, $g_0=Ah_0$, $g_m=Ah_m$, and
$\widehat g=A\widehat h_n$.  By the defining inequality,
\[
-2Y(\widehat g)+\|\widehat g\|_2^2
\le
-2Y(g_m)+\|g_m\|_2^2.
\]
Substituting $Y(g)=\langle g_0,g\rangle+n^{-1/2}\mathbb W(g)$ gives
\[
\|\widehat g-g_0\|_2^2
\le
\|g_m-g_0\|_2^2
+
\frac2{\sqrt n}\mathbb W(\widehat g-g_m).
\]
The set $A\cH_m-Ag_m$ lies in a linear space of dimension
$M=\dim V_m\asymp m^d$ and has bounded diameter.  The standard
finite-dimensional Gaussian maximal inequality therefore gives
\[
\E\sup_{e\in A\cH_m-Ag_m}|\mathbb W(e)|
\le
C\sqrt{M}\log^{1/2}(2+m).
\]
Consequently
\[
\E\|\widehat g-g_m\|_2^2
\le
C\left(
\|g_m-g_0\|_2^2+\frac{M\log m}{n}
\right).
\]
Using \eqref{eq:P-L2}, \eqref{eq:P-forward}, and the definition of $m$,
\[
\E\|\widehat h_n-h_0\|_2^2
\le
2\|h_m-h_0\|_2^2
+
2\E\|\widehat h_n-h_m\|_2^2.
\]
By \eqref{eq:restricted-lower},
\[
\|\widehat h_n-h_m\|_2
\le
Cm\|\widehat g-g_m\|_2.
\]
Therefore
\[
\E\|\widehat h_n-h_0\|_2^2
\le
C\left(
m^{-2\alpha}
+
m^2\frac{m^d\log m}{n}
\right)
\le
Cn^{-\frac{2\alpha}{2\alpha+d+2}}\log n.
\]

\emph{Lower bound.}
Let $M\asymp m^d$ be the number of top-scale good wave packets in
Proposition~\ref{prop:microlocal}.  Let $\psi_1,\dots,\psi_M$ be their
$L^2$-normalized representatives.  They satisfy
\begin{equation}\label{eq:packet-bounds}
\|\psi_j\|_{C^\alpha}\le Cm^{\alpha+d/2},
\qquad
\|A\psi_j\|_2\ge cm^{-1},
\end{equation}
and their forward images are asymptotically orthogonal at scale $m^{-1}$.
Set
\[
b=\eta r m^{-\alpha-d/2}
\]
with $\eta>0$ sufficiently small.  By Varshamov--Gilbert, choose
$\mathcal E\subset\{-1,1\}^M$ such that
\[
\log|\mathcal E|\ge cM
\]
and every distinct $\varepsilon,\varepsilon'\in\mathcal E$ have Hamming
distance at least $cM$.  Define
\[
h_\varepsilon=b\sum_{j=1}^M\varepsilon_j\psi_j.
\]
Finite-overlap wave-packet estimates give
\[
\|h_\varepsilon\|_{C^\alpha}\le r.
\]
If $\varepsilon,\varepsilon'$ differ in $D(\varepsilon,\varepsilon')$
coordinates, then
\[
\|h_\varepsilon-h_{\varepsilon'}\|_2^2
\asymp
b^2D(\varepsilon,\varepsilon'),
\]
and
\[
\|Ah_\varepsilon-Ah_{\varepsilon'}\|_2^2
\asymp
b^2m^{-2}D(\varepsilon,\varepsilon').
\]
For the Gaussian measures
$\Prob_\varepsilon=N(Ah_\varepsilon,n^{-1}I)$,
\[
\KL(\Prob_\varepsilon\|\Prob_{\varepsilon'})
=
\frac n2\|Ah_\varepsilon-Ah_{\varepsilon'}\|_2^2.
\]
At the choice \eqref{eq:balance},
\[
n b^2m^{-2}
=
\eta^2r^2
n m^{-2\alpha-d-2}
=O(1).
\]
Thus the average KL divergence is $O(M)$, while
$\log|\mathcal E|\asymp M$.  Fano's inequality gives a constant lower bound
for the probability of recovering a constant fraction of the coordinates
incorrectly.  Therefore
\[
\inf_{\widehat h}\sup_\varepsilon
\E_\varepsilon\|\widehat h-h_\varepsilon\|_2^2
\gtrsim
Mb^2
\asymp
m^{-2\alpha}.
\]
This proves the linear lower bound.
\end{proof}

\section{Lower bound for the nonlinear problem}

\begin{theorem}\label{thm:lower}
For $r$ sufficiently small,
\[
\inf_{\widehat f_n}
\sup_{f_0\in\cF(\alpha,r)}
\E_{f_0}\|\widehat f_n-f_0\|_2^2
\ge
c n^{-\frac{2\alpha}{2\alpha+d+2}}.
\]
\end{theorem}

\begin{proof}
Use the hypercube $\{h_\varepsilon:\varepsilon\in\mathcal E\}$ from the lower
bound in Theorem~\ref{thm:linear}, with
\[
m\asymp n^{1/(2\alpha+d+2)}.
\]
Set
\[
f_\varepsilon=1+h_\varepsilon.
\]
By construction, $f_\varepsilon\in\cF(\alpha,r)$.

For $\varepsilon,\varepsilon'$, let
\[
\delta=h_\varepsilon-h_{\varepsilon'}.
\]
From the wave-packet construction,
\[
\|\delta\|_{L^\infty}
\le
C r m^{-\alpha},
\qquad
\|A\delta\|_2
\ge
cm^{-1}\|\delta\|_2.
\]
Lemma~\ref{lem:quadratic}, applied with one center equal to $0$, gives
\[
\|F(h_\varepsilon)-F(h_{\varepsilon'})-A\delta\|_2
\le
C\|\delta\|_\infty\|\delta\|_2
\le
Crm^{-\alpha}\|\delta\|_2.
\]
Consequently
\[
\frac{\|F(h_\varepsilon)-F(h_{\varepsilon'})\|_2}
{\|A\delta\|_2}
\ge
1-Crm^{1-\alpha}.
\]
Since $\alpha>1$, this tends to $1$ as $m\to\infty$.  Hence, for all large
$m$,
\begin{equation}\label{eq:nonlinear-distance}
\|u_{f_\varepsilon}-u_{f_{\varepsilon'}}\|_2^2
\ge
c\,b^2m^{-2}D(\varepsilon,\varepsilon').
\end{equation}
The corresponding white-noise laws satisfy
\[
\KL(\Prob_{f_\varepsilon}\|\Prob_{f_{\varepsilon'}})
=
\frac n2\|u_{f_\varepsilon}-u_{f_{\varepsilon'}}\|_2^2.
\]
Exactly as in the linear case, the average KL divergence is $O(M)$ and
Fano's inequality yields
\[
\inf_{\widehat f}\sup_\varepsilon
\E_{f_\varepsilon}\|\widehat f-f_\varepsilon\|_2^2
\gtrsim
Mb^2
\asymp
m^{-2\alpha}.
\]
The right-hand side is the desired lower bound.
\end{proof}

\section{Upper bound: a local sieve estimator}

For $g\in L^2(D)$ define the formal least-squares functional
\[
\ell_n(g)=-2Y_n(g)+\|g\|_2^2.
\]
This notation is understood through any orthonormal expansion of $Y_n$; only
finite-dimensional projections are needed.

Choose
\[
m=m_n=\left\lfloor
n^{\frac1{2\alpha+d+2}}
\right\rfloor
\]
and set
\[
\cH_m
=
\left\{h\in V_m:\|h\|_{C^\alpha}\le r\right\}.
\]
Define
\begin{equation}\label{eq:sieve-estimator}
\widehat h_n
=
\arg\min_{h\in\cH_m}
\ell_n\bigl(u_{1+h}\bigr),
\qquad
\widehat f_n=1+\widehat h_n.
\end{equation}

\begin{theorem}\label{thm:upper}
For $r$ sufficiently small,
\[
\sup_{f_0\in\cF(\alpha,r)}
\E_{f_0}\|\widehat f_n-f_0\|_2^2
\le
C n^{-\frac{2\alpha}{2\alpha+d+2}}(\log n)^q
\]
for some finite $q$.
\end{theorem}

\begin{proof}
Write $f_0=1+h_0$ and $h_m=P_mh_0$.  By \eqref{eq:P-stable}, after replacing
$h_m$ by a fixed contraction if necessary, $h_m\in\cH_m$.  Put
\[
g_0=u_{1+h_0},\qquad g_m=u_{1+h_m},\qquad
\widehat g=u_{1+\widehat h_n}.
\]

First,
\[
\|g_m-g_0\|_2
\le
\|A_{1+h_0}(h_m-h_0)\|_2
+
C\|h_m-h_0\|_\infty\|h_m-h_0\|_2
\]
by Lemma~\ref{lem:quadratic}.  The approximation estimates
\eqref{eq:P-L2}--\eqref{eq:P-forward} and the usual high-frequency
$L^\infty$ Jackson estimate give
\begin{equation}\label{eq:approx-forward}
\|g_m-g_0\|_2\le C m^{-\alpha-1}.
\end{equation}

The least-squares inequality gives
\[
\|\widehat g-g_0\|_2^2
\le
\|g_m-g_0\|_2^2
+
\frac2{\sqrt n}\mathbb W(\widehat g-g_m).
\]
Because $\widehat g-g_m$ ranges over a finite-dimensional set parametrized by
$M\asymp m^d$ variables, entropy and Gaussian chaining bounds give
\begin{equation}\label{eq:empirical-process}
\E\sup_{h\in\cH_m}
\left|
\mathbb W\bigl(u_{1+h}-u_{1+h_m}\bigr)
\right|
\le
C\sqrt{M\log m}.
\]
A standard quadratic peeling argument then yields
\[
\E\|\widehat g-g_m\|_2^2
\le
C\left(
m^{-2\alpha-2}+\frac{M\log m}{n}
\right).
\]
See \cite{vandeGeer} for the empirical-process inequality.

Now apply the nonlinear restricted stability estimate
\eqref{eq:nonlinear-restricted} to $h=\widehat h_n$ and $k=h_m$:
\[
\|\widehat h_n-h_m\|_2
\le
Cm\|\widehat g-g_m\|_2.
\]
Therefore
\[
\E\|\widehat h_n-h_m\|_2^2
\le
C\left(
m^{-2\alpha}
+
m^2\frac{m^d\log m}{n}
\right).
\]
Since $m^{2\alpha+d+2}\asymp n$, both terms have order
$m^{-2\alpha}\log m$.  Finally,
\[
\|\widehat f_n-f_0\|_2^2
=
\|\widehat h_n-h_0\|_2^2
\le
2\|\widehat h_n-h_m\|_2^2
+
2\|h_m-h_0\|_2^2,
\]
and \eqref{eq:P-L2} completes the proof.
\end{proof}

\begin{proof}[Proof of Theorem~\ref{thm:main}]
The lower bound is Theorem~\ref{thm:lower}.  The upper bound is
Theorem~\ref{thm:upper}.
\end{proof}

\section{Why the exponent has the displayed form}

The rate can be read directly from the order-one inverse estimate.  At spatial
resolution $m$:

\begin{enumerate}
\item A $C^\alpha$ perturbation can be approximated in $L^2$ with squared
bias $m^{-2\alpha}$.

\item The dimension of the resolution space is $M\asymp m^d$.

\item Since the derivative is one-smoothing on the good resolution space,
white-noise variance in the parameter scale is
\[
\frac{m^2M}{n}
\asymp
\frac{m^{d+2}}n.
\]

\item Balancing bias and variance,
\[
m^{-2\alpha}\asymp \frac{m^{d+2}}n,
\]
gives
\[
m^{2\alpha+d+2}\asymp n.
\]
\end{enumerate}

Thus the mean squared error is $m^{-2\alpha}$, namely
\[
n^{-\frac{2\alpha}{2\alpha+d+2}}.
\]

The smallness of $r$ is used to ensure that the quadratic part of the forward
map is dominated by the order-one Fr\'echet derivative on the relevant
microlocal scales.  The condition $\alpha>1$ is exactly what gives an extra
high-frequency gain in the quadratic term.

\begin{thebibliography}{99}

\bibitem{BirmanSolomyak}
M.~\v{S}. Birman and M.~Z. Solomyak,
\emph{Asymptotic behavior of the spectrum of pseudodifferential operators with
anisotropically homogeneous symbols},
Vestnik Leningrad Univ. Math. \textbf{10} (1982), 237--247.

\bibitem{Bony}
J.-M. Bony,
\emph{Calcul symbolique et propagation des singularit\'es pour les \'equations
aux d\'eriv\'ees partielles non lin\'eaires},
Ann. Sci. \'Ecole Norm. Sup. \textbf{14} (1981), 209--246.

\bibitem{Donoho}
D.~L. Donoho,
\emph{Nonlinear solution of linear inverse problems by wavelet-vaguelette
decomposition},
Ann. Statist. \textbf{23} (1995), 705--742.

\bibitem{HormanderIII}
L. H\"ormander,
\emph{The Analysis of Linear Partial Differential Operators III},
Springer, Berlin, 1985.

\bibitem{Nickl2020}
R. Nickl, S. van de Geer, and S. Wang,
\emph{Convergence rates for penalized least squares estimators in
PDE-constrained regression problems},
arXiv:1809.08818.
\href{https://export.arxiv.org/pdf/1809.08818v3.pdf}{export.arxiv.org}.

\bibitem{NicklBayes}
R. Nickl,
\emph{Consistency of Bayesian inference with Gaussian process priors in an
elliptic inverse problem},
Inverse Problems \textbf{36} (2020), 085011.
\href{https://google.iopscience.iop.org/article/10.1088/1361-6420/ab7d2a}{google.iopscience.iop.org}.

\bibitem{PDERegressionMAP}
\emph{Convergence rates for the MAP estimator in PDE-regression models with
random design},
arXiv:2409.03417.
\href{https://arxiv.org/html/2409.03417}{arxiv.org}.

\bibitem{GlobalPoly2026}
\emph{Global polynomial-time estimation in statistical nonlinear inverse
problems via generalized stability},
arXiv:2601.09007.
\href{https://arxiv.org/html/2601.09007}{arxiv.org}.

\bibitem{Tsybakov}
A.~B. Tsybakov,
\emph{Introduction to Nonparametric Estimation},
Springer, New York, 2009.

\bibitem{vandeGeer}
S. van de Geer,
\emph{Empirical Processes in M-Estimation},
Cambridge University Press, Cambridge, 2000.

\bibitem{Zworski}
M. Zworski,
\emph{Semiclassical Analysis},
American Mathematical Society, Providence, 2012.

\end{thebibliography}

% Web-search provenance:
% [google.iopscience.iop.org](https://google.iopscience.iop.org/article/10.1088/1361-6420/ab7d2a)
% [arxiv.org](https://arxiv.org/html/2409.03417)
% [export.arxiv.org](https://export.arxiv.org/pdf/1809.08818v3.pdf)
% [arxiv.org](https://arxiv.org/html/2601.09007)
% [exa.ai](https://exa.ai/library/publication/rj7jk1yjt53)

\end{document}
```